\section{Islands 5 and 7: moving targets, a distinguished seed, and geometric attachment}\label{sec:dynamics}
\subsection{Exact arithmetic encoding and a metric hitting theorem}
Let $\widehat\ZZ$ be the group of profinite integers, let $m_H$ be Haar probability measure, and let $T(\xi)=\xi+1$.
We write $m_H$ to avoid confusion with the M\"obius function $\mu$. Set
\[
y_n=\sqrt{n+2},\quad Q_n=\prod_{p\le y_n}p,\quad
E_n=\{\xi:p\nmid\xi(\xi+2)\text{ for every }p\le y_n\},\quad
Z_n(\xi)=\1_{E_n}(T^n\xi).
\]
For $n\ge3$, $Z_n(0)=1$ if and only if both $n$ and $n+2$ are prime: both exceed the common sieve threshold, whereas every composite integer at most $n+2$ has a prime factor at most $\sqrt{n+2}$.
The Chinese remainder theorem gives
\[
\alpha_n=m_H(E_n)=\frac12\prod_{2<p\le y_n}(1-2/p)\asymp(\log n)^{-2}.
\]
The underlying $T$ is an autonomous adding map, not a mixing transformation. Nonautonomy enters through the targets and sieve scales that change with $n$.

\begin{theorem}[Haar moving-target variance and almost-everywhere hitting]\label{thm:haar}
Let $S_N=\sum_{3\le n\le N}Z_n$ and $A_N=\sum_{3\le n\le N}\alpha_n$. Then
\[
\operatorname{Var}_{m_H}(S_N)\ll N,\qquad A_N\asymp N/(\log N)^2.
\]
For $m_H$-almost every $\xi$, $S_N(\xi)\to\infty$.
\end{theorem}
\begin{proof}
For $m<n$, put $d=n-m$. Separating common sieve primes from newly added primes in the Chinese remainder theorem gives
\[
\operatorname{Cov}(Z_m,Z_n)=\alpha_m\alpha_n(K_{y_m}(d)-1).
\]
For a common prime $p$, the forbidden classes are $0,-2,-d,-d-2$. Counting their coincidences gives
\begin{align*}
k_2(d)&=2\1_{2\mid d},&k_3(d)&=3\1_{3\mid d},\\
k_p(d)&=1-\frac4{(p-2)^2}+\frac{2p}{(p-2)^2}\1_{p\mid d}
+\frac p{(p-2)^2}(\1_{p\mid d-2}+\1_{p\mid d+2})\quad(p\ge5),
\end{align*}
where $K_y=\prod_{p\le y}k_p$. Every factor has mean $1$ over a complete residue system.
Expanding into indicators of Chinese-remainder residue classes, the sum of absolute coefficients is at most
\[
D(y)=6\prod_{5\le p\le y}(1-2/p)^{-2}\ll\log^4(2y).
\]
The final bound follows from Mertens' product theorem and a convergent $O(p^{-2})$ correction.
The count of any residue class in an integer interval differs from the interval length divided by its modulus by at most $1$. Hence
\[
\left|\sum_{d\in J}(K_y(d)-1)\right|\le D(y).
\]
For fixed $m$, Abel summation against the nonincreasing weights $\alpha_{m+d}$ yields
\[
\left|\sum_{d=1}^{N-m}\alpha_{m+d}(K_{y_m}(d)-1)\right|
\le\alpha_{m+1}D(y_m).
\]
When $y_m\ge3$, $\alpha_m^2D(y_m)=1/6$, so the absolute value of the grouped covariance for each $m$ is at most $1/6$.
The finitely many initial $m$ contribute a bounded amount, and the diagonal variance is at most $\sum\alpha_n\le N$. This proves the variance estimate.
Mertens' product theorem and the fixed two-point singular series give the order of $A_N$.
Chebyshev's inequality along $N=j^2$ gives
\[
m_H\{|S_{j^2}-A_{j^2}|>A_{j^2}/2\}\ll\frac{(\log j)^4}{j^2}.
\]
The right-hand side is summable. The first Borel--Cantelli lemma implies that almost every seed eventually satisfies $S_{j^2}\ge A_{j^2}/2\to\infty$. Monotonicity of $S_N$ extends the conclusion to all $N$.
\end{proof}
This argument comes from the original moving-target proof\cite{archivedynamics}. It is a metric theorem for the actual sieve events, not merely a finite-model statement.
Its null exceptional set may nevertheless contain $0$, so it does not cross the map's Bridge B.

\subsection{A precise form of the distinguished-seed gap}
The following sufficient criterion comes from the existing distinguished-seed reconnaissance report\cite{archiveseed}; we reproduce its short proof.
Choose $q_N\mid Q_N$, put $U_N=q_N\widehat\ZZ$, and define
\[
V_N=\|S_N-A_N\|_{L^2(m_H)},\qquad
\Omega_N=\sup_{\xi\in U_N}(S_N(\xi)-S_N(0)).
\]
\begin{proposition}[A sufficient cylinder-set return criterion]\label{prop:seed}
For every $N$,
\[
S_N(0)\ge A_N-\sqrt{q_N}V_N-\Omega_N.
\]
If, along an unbounded sequence $N_j$, some fixed $\eta>0$ satisfies
$\Omega_{N_j}+\sqrt{q_{N_j}}V_{N_j}\le(1-\eta)A_{N_j}$, then there are infinitely many twin primes.
\end{proposition}
\begin{proof}
We have $m_H(U_N)=1/q_N$. The definition of $\Omega_N$ and Cauchy--Schwarz imply
\[
S_N(0)\ge q_N\int_{U_N}S_N\,dm_H-\Omega_N
\ge A_N-\sqrt{q_N}V_N-\Omega_N.
\]
The assumed condition gives $S_{N_j}(0)\ge\eta A_{N_j}\to\infty$; apply the exact arithmetic encoding.
\end{proof}
The known bound $V_N\ll\sqrt N$ ensures that $q_N=o(N/\log^4N)$ implies
$\sqrt{q_N}V_N=o(A_N)$. This is sufficient, not necessary, and the required bound for $\Omega_N$ is not known in this regime.
For the fine cylinder $q_N=Q_N$, $S_N$ is constant on the cylinder and $\Omega_N=0$, but
$Q_N=\exp((1+o(1))\sqrt N)$ makes the generic point-evaluation cost too large.
The proposition is a sufficient criterion, not an assertion that its hypothesis holds or that it is the only possible return mechanism.

\subsection{Stage-preserving sieve evolution and a scoped obstruction to strongly mixing carriers}
Let $P_k=\prod_{j\le k}p_j$, equip $V_k=\CC^{\ZZ/P_k\ZZ}$ with the normalized counting inner product, and put
$B_k(r)=\1_{(r(r+2),P_k)=1}$, $a_k=\operatorname{mean}B_k$, and $W_k=B_k-a_k$.
When a new odd prime $p$ is added, define the replication-and-deletion operator by
\[
(R_pf)(r+jP_k)=f(r)\1_{p\nmid(r+jP_k)(r+jP_k+2)},\quad 0\le j<p.
\]
Exactly $p-2$ children survive for every parent residue. Thus, with $\alpha_p=1-2/p$,
\[
R_p^*R_p=\alpha_p I,\quad \operatorname{mean}(R_pf)=\alpha_p\operatorname{mean}f.
\]
Set $g=R_p1-\alpha_p$. Direct averaging gives $\|g\|^2=\alpha_p(1-\alpha_p)$, and every mean-zero $f$ satisfies $\langle R_pf,g\rangle=0$. Consequently,
\[
W_{k+1}=R_pW_k+a_kg,\qquad \|W_k/a_k\|^2=a_k^{-1}-1\asymp(\log p_k)^2.
\]
The forcing term is part of the actual evolution; this is not a homogeneous exponentially contracting system.

More precisely, suppose maps $J_k$ are uniformly bounded and coercive on the entire mean-zero space,
$c\|f\|\le\|J_kf\|\le C\|f\|$, and satisfy the exact intertwining relation
$J_{k+1}R_{p_{k+1}}=Q_kJ_k$.
They cannot also satisfy a uniform arbitrary-product estimate
$\|Q_{b-1}\cdots Q_a\|\le C_Q\theta^{b-a}$ with $0<\theta<1$.
Indeed, the three assumptions imply
\[
c\prod_{a<j\le b}(1-2/p_j)^{1/2}\le CC_Q\theta^{b-a}.
\]
For fixed $a$, the left side is only of order $1/\log p_b$, and the prime number theorem gives $p_b\asymp b\log b$, contradicting exponential decay on the right.
This excludes only the simultaneous validity of those three assumptions. It does not exclude noncoercive quotient maps, observable seminorms, or nonautonomous representations with controlled defects\cite{archiveintertwiner}.

\subsection{The role and limits of Island 7}
A useful geometric lift must at least supply a time-preserving semiconjugacy, event pullback, an applicable measure correspondence, and a correspondence for the \emph{distinguished seed}.
For example, candidate maps $\rho_k$ and transformations $H_k,F_k$ must satisfy
$\rho_{k+1}H_k=F_k\rho_k$, and lifted events must be exactly the inverse images of the original events.
These are attachment conditions to be proved. We do not construct a twin-prime-specific H\'enon lift here.

Even a stationary mixing system cannot have a nontrivial addition-by-one clock modulo $M$ as a measure-preserving factor: the clock's mean-zero eigenfunction $\e(r/M)$ has autocorrelation of constant absolute value $1$, contradicting mixing after pullback.
This finite-factor argument does not establish impossibility for all nonautonomous geometric representations.
Nor does a two-dimensional phase space imply that all relevant observables have a two-dimensional linear carrier.
Island 7 is therefore an auxiliary representation route; a geometric picture cannot replace the distinguished-orbit hitting requirement on Island 5.
